\section{Finite clopen routing}
\label{sec:routing}

Throughout this section, let

\begin{equation}
  X=\bigsqcup_{i\in I}X_i,
  \label{eq:finite-clopen-partition}
\end{equation}

where $I$ is finite and every $X_i$ is nonempty, clopen, invariant, and
irreducible.  Pairwise disjointness and compactness provide the separation
that preinjectivity needs.

\begin{lemma}[Homoclinic separation]
\label{lem:homoclinic-separation}
Every homoclinicity class in $X$ is contained in one $X_i$.  Consequently, for
an endomorphism $T:X\to X$,
\begin{equation}
 T\text{ is preinjective}
 \quad\Longleftrightarrow\quad
 T|_{X_i}\text{ is preinjective for every }i\in I.
 \label{eq:componentwise-pi}
\end{equation}
The equivalence does not require the sets $T(X_i)$ to be disjoint.
\end{lemma}

\begin{proof}
Fix $i\ne j$.  The disjoint compact sets $X_i$ and $X_j$ have positive
distance $\delta_{ij}$ in a compatible product metric.  If $x\in X_i$ and
$y\in X_j$, invariance gives
\[
 d(\sigma^n x,\sigma^n y)\geq\delta_{ij}
 \qquad(n\in\Z).
\]
If $x\sim_h y$, their finite disagreement set escapes every central
coordinate window under $\sigma^n$ as $n\to+\infty$ and as $n\to-\infty$.
Thus $d(\sigma^n x,\sigma^n y)$ tends to zero in both tails, a contradiction.
Every homoclinic pair therefore lies in one piece.

The forward implication in \eqref{eq:componentwise-pi} follows by
restriction.  Conversely, every homoclinic collision for $T$ has both sources
in one $X_i$ and is excluded by the corresponding restriction.  A collision
between sources in different pieces has no bearing on preinjectivity.
\end{proof}

\begin{lemma}[Finite clopen pasting]
\label{lem:clopen-pasting}
Suppose $T_i:X_i\to X$ is continuous and shift commuting for every $i\in I$.
The piecewise map $T:X\to X$ defined by $T|_{X_i}=T_i$ is a cellular
automaton.
\end{lemma}

\begin{proof}
A map continuous on every member of a finite clopen partition is continuous
on the union.  Because each $X_i$ is invariant, $T\sigma=\sigma T$ can be
checked on each piece.  Hence $T$ is continuous and equivariant.  The
Curtis--Hedlund characterization supplies a finite memory and one local rule
on all of $X$.
\end{proof}

The target of a piece is also well defined.

\begin{lemma}[Unique component image]
\label{lem:unique-component-image}
If $T:X\to X$ is an endomorphism, then for every $i\in I$ there is a unique
$\rho_T(i)\in I$ such that
\begin{equation}
  T(X_i)\subset X_{\rho_T(i)}.
  \label{eq:component-route}
\end{equation}
\end{lemma}

\begin{proof}
The restriction $T|_{X_i}$ maps the irreducible system $X_i$ onto its image,
so $T(X_i)$ is irreducible.  Suppose it met $X_j$ and $X_k$ with $j\ne k$.
The intersections with those invariant clopen sets would be nonempty
relatively open subsets of $T(X_i)$, and no shift iterate of one could meet
the other.  This contradicts irreducibility.  The image is nonempty, so the
target index exists and is unique.
\end{proof}

The preceding lemmas close the routing criterion.

\begin{theorem}[Finite-clopen routing criterion]
\label{thm:routing}
For the decomposition \eqref{eq:finite-clopen-partition},
\begin{equation}
 X\text{ is Myhill}
 \quad\Longleftrightarrow\quad
 \begin{cases}
 X_i\text{ is Myhill for every }i\in I,\\
 X_i\not\pito X_j\text{ for all }i\ne j.
 \end{cases}
 \label{eq:routing-theorem}
\end{equation}
\end{theorem}

\begin{proof}
Suppose some $X_i$ is not Myhill.  Choose a preinjective nonsurjective
endomorphism $S_i:X_i\to X_i$ and paste it with the identity on every other
piece.  By \cref{lem:clopen-pasting} this is an endomorphism of $X$, and by
\cref{lem:homoclinic-separation} it is preinjective.  A point in
$X_i\setminus S_i(X_i)$ cannot be supplied by an unchanged disjoint piece, so
the pasted map is not onto.  Thus $X$ is not Myhill.

Now suppose $i\ne j$ and let $\phi:X_i\pito X_j$.  Define
\begin{equation}
 T(x)=
 \begin{cases}
 \phi(x),&x\in X_i,\\
 x,&x\in X_k,\ k\ne i.
 \end{cases}
 \label{eq:redirect}
\end{equation}
Again \cref{lem:clopen-pasting,lem:homoclinic-separation} make $T$ a
preinjective endomorphism.  Its image may contain collisions between
$\phi(X_i)$ and the identity image of $X_j$, but their sources lie in
different homoclinicity classes.  More importantly,
\[
 T(X_i)\subset X_j,
 \qquad T(X_k)=X_k\quad(k\ne i),
\]
so no point of the source component $X_i$ belongs to $T(X)$.  Hence $T$ is
not onto and $X$ is not Myhill.

Conversely, assume both conditions on the right of
\eqref{eq:routing-theorem}, and let $T:X\to X$ be preinjective.  By
\cref{lem:unique-component-image}, its restriction gives a preinjective arrow
\[
 T|_{X_i}:X_i\pito X_{\rho_T(i)}.
\]
The absence of off-diagonal arrows forces $\rho_T(i)=i$.  Since $X_i$ is
Myhill, the restriction maps $X_i$ onto itself.  Taking the union over the
finite partition gives $T(X)=X$.
\end{proof}

\begin{remark}[Why the hypotheses are visible]
\label{rem:routing-scope}
An infinite union can fail the continuity step at an accumulation set.  A
general reducible SFT can contain heteroclinic or wandering configurations
outside its irreducible components, so the displayed pieces need not exhaust
the homoclinicity classes or support a piecewise definition on the whole
space.  Neither situation is covered by \cref{thm:routing}.
\end{remark}

